\clearpage
\subsection{Flag Meanings}
Integer condition codes are held in FLAGS. Integer instructions leave FLAGS unchanged unless their instruction
description explicitly defines a FLAGS write. Every FLAGS-writing instruction replaces the complete Z, N, C, and V
image; partial FLAGS writes do not exist. An instruction that does not write FLAGS preserves the complete image.

For an operand width of n bits, ordinary integer ALU results are evaluated modulo 2\textasciicircum{}n unless the
instruction explicitly defines a wider intermediate.

Explicit FLAGS-writing integer instructions include CMP, TEST, ADC, SBB, INCF, DECF, BTEST, BSET, BCLR, BCHG,
SETF, WRFLAGS, CMPXCHG, SEGLEA, and the bounds-check instructions. CMPJcc and TESTJcc compute temporary condition flags
for their branch decision and do not update architectural FLAGS. ADD, SUB, logic operations, INC, DEC, shifts,
rotates, moves, extensions, min/max, and multiply/divide instructions leave FLAGS unchanged.

\manualtablecaption{Integer Condition-Code Bits}
\begin{manualtabular}{@{}>{\raggedright\arraybackslash}p{0.45in}>{\raggedright\arraybackslash}p{1.05in}>{\raggedright\arraybackslash}p{3.90in}@{}}
\toprule
\rowcolor{ManualHeaderFill}
\textbf{Bit} & \textbf{Name} & \textbf{Meaning}\\
\midrule
Z & zero & Set when the result value is zero.\\
N & negative & Set from the most significant result bit at the operand size.\\
C & carry/borrow & Set on unsigned carry out for addition or unsigned borrow for subtraction.\\
V & overflow & Set on signed overflow or an explicitly reported exceptional condition.\\
\bottomrule
\end{manualtabular}

\subsection{Conditional Consumers}
Conditional branches, calls, traps, and sets consume the condition-code table without recomputing FLAGS. An ordinary
conditional consumer observes the FLAGS image produced by the last committed FLAGS-writing instruction. REPcc instead
tests a temporary flag image derived from the body instruction\textquotesingle{}s repeat-observed value, as defined in Repeated Scalar
Execution. That test leaves architectural FLAGS unchanged.

A false condition suppresses the conditional action unless the instruction description explicitly defines a different
commit rule, such as a repeat instruction\textquotesingle{}s terminating iteration rule.

\subsection{Floating-Point Interactions}
Floating-point comparison instructions may update integer condition codes only when their instruction definition says
so. FFLAGS holds floating-point exception flags, and FSTATUS holds the corresponding exception-condition enables
and rounding state. IEEE-754 FFLAGS causes use independent accrued semantics. The complete-image rule applies to
integer FLAGS.

Integer conditional consumers do not read FFLAGS directly; floating-point instructions that want integer control-flow
consumers must materialize the selected condition into FLAGS or an Rn value.
